% Validator Design Rationale Subsection
% For inclusion in paper Section III-C (Deterministic Validation Pipeline)

\subsection{Deterministic Validation Pipeline}
\label{sec:validator}

A key design principle of our system is that the LLM compiler operates
\textit{offline} and its output is never trusted directly. Every
ConstraintProgram passes through an 8-pass deterministic validator before
reaching the routing layer. This design is motivated by two observations:
(1)~LLMs can produce syntactically valid but semantically incorrect
constraint programs, and (2)~in safety-critical network infrastructure,
a single undetected constraint violation can cascade into service outages.

The validator implements the following passes in sequence, with early
termination on fatal errors:

\begin{enumerate}
\item \textbf{Schema Validation.} Checks structural completeness: all
required fields present, valid priority levels, non-empty constraint
types and targets. Rejects malformed programs before deeper analysis.

\item \textbf{Entity Grounding.} Verifies that all referenced entities
exist in the constellation model: node IDs $\in [0, N)$, plane IDs
$\in [0, P)$, region names $\in \mathcal{R}$, traffic classes
$\in \mathcal{T}$. This catches hallucinated entities (e.g., node~454
in a 400-node constellation).

\item \textbf{Type Safety.} Ensures constraints attach to semantically
correct entity types: \texttt{max\_latency\_ms} must target a
\texttt{flow\_selector}, \texttt{disable\_node} must target a
\texttt{node}, \texttt{avoid\_latitude} must target \texttt{edges}.
Prevents type confusion errors where the LLM assigns a constraint
to the wrong entity class.

\item \textbf{Value Range Checking.} Validates numeric parameters:
latitudes $\in [-90, 90]$, latency values $> 0$, utilization caps
$\in (0, 1]$, node IDs within constellation bounds. Catches
out-of-range values that would produce undefined behavior.

\item \textbf{Conflict Detection.} Identifies contradictory constraints
within the same program: a node cannot be simultaneously disabled and
used as a routing waypoint; conflicting latency bounds on the same flow
are flagged. Contradictions are promoted to errors (not warnings),
ensuring logically inconsistent programs are rejected.

\item \textbf{Physical Admissibility.} Checks whether the constrained
topology is physically realizable: latency deadlines below the
single-hop physical minimum ($<2.0$\,ms) are rejected; latitude
avoidance thresholds that remove $>50\%$ of edges trigger warnings.

\item \textbf{Reachability Analysis.} Performs BFS on the constrained
graph $G' = (V \odot M_V, E \odot M_E)$ to verify connectivity.
Severe capacity loss ($\geq 75\%$ of nodes disabled) triggers a strong
warning; moderate loss ($\geq 50\%$) triggers a standard warning.
These capacity thresholds are heuristic indicators outside the
soundness path---they do not block acceptance, which is determined
solely by Pass~8.

\item \textbf{Feasibility Certification.} For each demanded flow,
constructs a routing witness on the constrained topology to certify
that all hard constraints can be simultaneously satisfied. Five
certified fragments cover the constraint space:
\begin{itemize}
\item \textbf{F1} (topology only): BFS reachability
\item \textbf{F2} (+ latency): Dijkstra with deadline
\item \textbf{F3} (+ hops): BFS with hop limit
\item \textbf{F4} (+ latency + hops): hop-layered Dijkstra
\item \textbf{F5} (+ $k$-disjoint): Edmonds-Karp max-flow
\end{itemize}
Three outcomes: \decision{accept} (witness found), \decision{reject}
(no feasible routing exists), or \decision{abstain} (unsupported
constraint combination). Programs are rejected only on \decision{reject};
\decision{abstain} defers to Dijkstra fallback routing, preserving
safety without constructive certification.
\end{enumerate}

\begin{theorem}[Acceptance Soundness]
If the feasibility certifier accepts a constraint program $P$ with
witness $W$, then there exists a routing assignment satisfying all
hard constraints of $P$ on the constrained topology.
\end{theorem}

\begin{proof}[Proof sketch]
By case analysis over fragments F1--F5. Each fragment's algorithm
is a standard shortest-path or max-flow algorithm whose correctness
is well-established. The witness $W$ is the concrete path (or path set)
returned by the algorithm, which by construction satisfies the
topology constraints (disabled nodes/edges excluded from the search
graph), latency bounds (Dijkstra optimality), hop limits (BFS layering),
and disjointness requirements (augmenting path decomposition).
Unsupported combinations produce \decision{abstain}, never
\decision{accept}, so no false acceptance is possible within the
certified fragment space.
\end{proof}

\noindent\textbf{Design rationale.} We chose deterministic
validation over learned or probabilistic checking for three reasons:
\begin{itemize}
\item \textit{Completeness}: every structural error class is covered by
at least one pass, achieving 100\% detection on our corruption benchmark
(8 error types $\times$ 30 injections = 240 tests).
\item \textit{Transparency}: each rejection includes a human-readable
error message identifying the specific violation, enabling the repair
loop to provide targeted feedback to the LLM.
\item \textit{Soundness}: the feasibility certifier guarantees that
accepted programs have constructive routing witnesses, closing the
semantic gap between structural validity and routing feasibility.
\end{itemize}

\noindent\textbf{Repair loop integration.} When validation fails, the
error messages are fed back to the LLM as a repair prompt. The compiler
retries up to $k=3$ times, with each attempt receiving the previous
errors as context. In our 240-benchmark evaluation (uniform random edge delays), 97.9\% of intents
compile successfully, with 93.3\% succeeding on the first attempt.
